\section{Results: Explainability and Cross-dataset Generalization}

\subsection{Grad-CAM}

To interrogate whether the student's in-distribution accuracy rests on cytologically meaningful evidence rather than incidental background cues, we visualized class-discriminative attention with Grad-CAM \citep{selvaraju2017gradcam} applied to the final convolutional block of the EfficientNet-B0 student. Figure~\ref{fig:gradcam} presents one representative cell from each of the five SIPaKMeD classes. All five exemplars were classified correctly, and at high confidence (predicted-class probability between $0.97$ and $1.00$), so the visualizations report attention for decisions the model got right rather than rationalizations of errors. Across every panel the salient regions concentrate on the nucleus and the surrounding cytoplasmic body, the same structures a cytotechnologist inspects when assessing nuclear enlargement, hyperchromasia, perinuclear haloing, and keratinization. The heatmaps do not localize to slide background, staining gradients, or crop borders, which would have signalled shortcut learning of the kind that frequently inflates apparent performance on medical-image benchmarks \citep{geirhos2020shortcut}. We therefore read these attributions as evidence that the in-distribution decision surface is anchored to plausible cellular morphology. This interpretation carries an important caveat: Grad-CAM is a saliency-attribution method that reports \emph{where} gradient-weighted activations are strongest, not \emph{why} a clinical category was chosen. Spatial overlap with the nucleus is consistent with sound reasoning but does not constitute a faithfulness guarantee, and it says nothing about whether the learned features survive a change of imaging domain---the question we turn to next.

\subsection{Cross-dataset performance}

Against the $0.966$ accuracy achieved within SIPaKMeD, performance on three independent external Pap datasets collapses toward chance. Table~\ref{tab:cross} reports the binary Normal-versus-Abnormal results and Figure~\ref{fig:cmcross} the corresponding confusion matrices. On Herlev \citep{jantzen2005herlev} ($n=917$) the model reaches accuracy $0.639$, balanced accuracy $0.651$, sensitivity $0.625$, specificity $0.678$, and AUROC $0.728$. On HiCervix \citep{cai2024hicervix} ($n=4823$) accuracy falls to $0.605$ with balanced accuracy $0.562$, a markedly depressed sensitivity of $0.386$, specificity $0.738$, and AUROC $0.594$. On APACC ($n=1125$) the model is effectively at chance: accuracy $0.532$, balanced accuracy $0.535$, sensitivity $0.654$, specificity $0.417$, and AUROC $0.525$. The internal-to-external degradation is therefore on the order of $33$ accuracy points in the most favourable case and larger elsewhere, and the discriminative ranking signal weakens monotonically from Herlev through HiCervix to APACC. This is the answer to RQ2 and confirms H2: strong, well-calibrated in-distribution accuracy does not, on its own, translate into cross-dataset generalization.

\subsection{Why cross-dataset evaluation fails}

We disentangle four contributing factors, and crucially distinguish a genuine model and domain failure from a mere artefact of our label crosswalk. \emph{First, the collapse is not a labelling artefact.} The class-conditional separation of the predicted abnormality score, $\Delta = \overline{P(\text{abnormal})}_{\text{abnormal}} - \overline{P(\text{abnormal})}_{\text{normal}}$, is positive on every dataset---$+0.284$ on Herlev, $+0.112$ on HiCervix, and $+0.045$ on APACC (Figure~\ref{fig:sep})---and all three AUROC values exceed $0.5$. Had our many-to-one mapping been inverted or ill-defined, the model would have ranked abnormal cells \emph{below} normal ones, yielding $\Delta<0$ and AUROC$<0.5$; instead the ordering is preserved, weakly but consistently, in the direction the mapping intends. Combined with an explicit audit of each mapper, this rules out relabelling as the cause (RQ3, H3).

\emph{Second, the dominant failure is operating-point miscalibration.} The in-distribution calibration (ECE $0.0152$) does not transfer: the baseline level of $P(\text{abnormal})$ shifts per domain, sitting high (around $0.6$) on APACC so the fixed threshold $\tau=0.5$ over-calls abnormality, and low (around $0.35$) on HiCervix so it under-calls it. A single global threshold consequently discards the ranking information the AUROC shows is present. The symptom is diagnostic: at $\tau=0.5$ accuracy falls \emph{below} the majority-class baseline for Herlev ($0.639$ vs.\ $0.74$) and HiCervix ($0.605$ vs.\ $0.62$), even though AUROC $>0.5$. Re-tuning $\tau$ on the test labels as an oracle recovers $0.75$ on Herlev (at $\tau=0.05$), $0.65$ on HiCervix (at $\tau=0.85$), and $0.53$ on APACC (at $\tau=0.40$)---confirming that much of the deficit is threshold placement rather than absent signal, and answering RQ4 (H4): calibration holds in-distribution but does not transfer.

\emph{Third, the residual signal degrades with slide-preparation distance from SIPaKMeD's conventional smear.} Threshold-free AUROC tracks this gradient cleanly: Herlev (conventional Pap) $0.728$ $>$ HiCervix (liquid-based cytology) $0.594$ $>$ APACC (auto-segmented noisy crops) $0.525\approx$ chance, indicating that preparation and segmentation domain shift, not labels, erode the discriminative features. \emph{Fourth, the internal $0.966$ is itself optimistic.} It derives from a random per-cell split that can place cells from the same cluster field-of-view on both sides of the train/validation boundary; a leakage-safe, cluster-grouped split would lower the internal figure, so the true internal-to-external gap is larger still than the headline numbers suggest. Together these factors recast the cross-dataset collapse as a problem of domain shift and operating-point transfer---addressable by stain and domain normalization \citep{macenko2009stain}, per-domain calibration \citep{guo2017calibration}, leakage-aware evaluation, and threshold-free reporting---rather than as a flaw in the label alignment.